% =====================================================================
%  Blaze2D — paper body (Design Choices … From Solver to Pipeline)
% =====================================================================

\section{Design Choices}\label{sec:design}

Three foundational choices were made from the start, each driven by a
concrete requirement of the envelope pipeline.

\paragraph{Method: plane-wave expansion.}
The PWE method discretizes the Bloch eigenproblem in reciprocal space
by expanding both the dielectric function and the field in Fourier
modes on the reciprocal lattice. For a given $\bm{k}$-point and
resolution (number of retained plane waves $N_{\mathrm{pw}}$), the
problem becomes a dense $N_{\mathrm{pw}} \times N_{\mathrm{pw}}$ matrix
eigenproblem. The reason for choosing PWE over real-space methods is
direct: the envelope framework requires not just eigenvalues but Bloch
functions and their Fourier coefficients, and the PWE basis delivers
these natively. Additionally, PWE eigenvalues converge spectrally with
resolution, providing high accuracy at moderate grid sizes~\cite{mpb}.

\paragraph{Algorithm: LOBPCG.}
The locally optimal block preconditioned conjugate gradient (LOBPCG)
method~\cite{knyazev2001} solves for the lowest eigenpairs of a
Hermitian operator without forming or factorizing the full matrix. It
is iterative, works in blocks (solving for several eigenpairs
simultaneously), and accepts a preconditioner. These properties make it
the natural choice for the photonic PWE problem, where only the lowest
bands are needed and a good preconditioner is available. MPB uses the
same algorithmic family~\cite{mpb}.

\paragraph{Language: Rust.}
The solver must run millions of operator applications across thousands
of crystal configurations. Rust provides predictable memory layout, no
garbage-collection pauses, and safe concurrency. These properties
matter once the solver runs inside a tight parameter-sweep loop. The
compiled performance is comparable to C/Fortran, while the type system
eliminates entire classes of memory bugs that would otherwise surface
only under heavy workloads.

\section{Geometry and the Permittivity Matrix}\label{sec:geometry}

The first module to build is the \emph{geometry}: translating a crystal
specification (lattice vectors, shapes, dielectric constants) into the
Fourier-space permittivity matrix $\epsilon_{\bm{G},\bm{G}'}$ processed
by the PWE operators.

Blaze2D uses the same analytical subpixel smoothing as MPB. Rather than
sampling $\epsilon(\bm{x})$ on a pixel grid, each geometric primitive
(cylinder, ellipse, block) is treated as an exact shape whose overlap
with every grid cell is computed analytically. The smoothed
$\epsilon^{-1}$ is then obtained via the anisotropic averaging
procedure of Farjadpour \textit{et~al.}~\cite{farjadpour2006}, which
avoids the staircasing artifacts that naive pixelization introduces at
dielectric interfaces.

Validating this module is straightforward: the permittivity matrix can
be compared element-by-element against MPB for any crystal geometry.
\Cref{fig:material_map_comparison} shows such a comparison for a square
lattice of dielectric rods ($\epsilon = 8.9$, $r = 0.45a$) at
resolutions $N = 8$ to $128$. Only the coarsest grid ($N = 8$) exhibits
appreciable error (mean $|\Delta\epsilon| \approx 3 \times 10^{-2}$,
peak $\approx 0.68$); for $N \geq 16$ the mean error drops below
$8 \times 10^{-3}$ and reaches $\approx 5 \times 10^{-4}$ at $N = 128$.
The remaining residual peaks of order $0.1$--$0.2$ are localized at the
rod--air interface, where the two solvers discretize the material
boundary slightly differently. Agreement to this level at every
resolution confirms that the geometry pipeline is correct before any
eigensolve is attempted.

\begin{figure}[t]
  \centering
  \includegraphics[width=\linewidth]{material_map_comparison}
  \caption{Smoothed permittivity maps from Blaze2D and MPB at
    resolutions $N = 8$ to $128$, with the bottom row showing the
    absolute difference $|\Delta\epsilon|$ on a power-law color scale
    ($\gamma = 0.3$) to magnify small discrepancies.}
  \label{fig:material_map_comparison}
\end{figure}

With the geometry validated, the remaining task is a spectral analysis:
construct the TE and TM operators on a Fourier basis and solve them
efficiently.

\section{Eigenvalue Operators and Their Application}\label{sec:operators}

With the permittivity matrix in hand, the next step is to discretize
the TE and TM master equations of the photonic-crystal eigenproblem
explicitly. In contrast to MPB, which works from the full curl--curl
operator and projects onto the relevant polarization, Blaze2D
implements the TE and TM operators directly:
\begin{alignat}{7}
  \text{TE:}\quad & A_{\mathrm{TE}} &{}={}& -(\bm{k} + \bm{G}) \cdot
  [\epsilon^{-1}]_{\bm{G},\bm{G}'}\,(\bm{k} + \bm{G}'),
  \qquad & B_{\mathrm{TE}} &{}={}& \identity, \label{eq:op_te}\\
  \text{TM:}\quad & A_{\mathrm{TM}} &{}={}& |\bm{k} + \bm{G}|^{2}\,\delta_{\bm{G},\bm{G}'},
  \qquad & B_{\mathrm{TM}} &{}={}& \epsilon_{\bm{G},\bm{G}'}. \label{eq:op_tm}
\end{alignat}
The operator classes are implemented generically: the LOBPCG solver
receives an abstract operator interface, making the TE and TM cases
interchangeable at the API level.

Operator application proceeds in the standard PWE fashion.
Multiplication by $\epsilon^{-1}$ in reciprocal space is a dense
matrix--vector product, while the gradient terms $(\bm{k} + \bm{G})$
act as diagonal multiplications. In practice, these operations dominate
the runtime and their memory-access pattern determines performance.

At high dielectric contrast, the Fourier coefficients of $\epsilon^{-1}$
decay slowly, and individual reciprocal-space components can develop
large magnitudes that amplify numerical noise. Blaze2D applies the same
clamping and regularization strategy as MPB to suppress these artifacts
without distorting the physically relevant low-frequency content.
This direct TE/TM implementation is also what makes the data extraction
convenient: the same operator objects used in the eigensolve can later
be reused to expose explicit matrix elements to the envelope pipeline.

\section{Preconditioning}\label{sec:preconditioning}

The convergence rate of LOBPCG depends critically on the
preconditioner. Each polarization admits a different strategy.

For TE, the preconditioner is the one described in Ref.~\cite{mpb}: an
approximation to $A_{\mathrm{TE}}^{-1}$ built from the diagonal kinetic
term $(|\bm{k} + \bm{G}|^{2})$ and a scalar estimate of
$\epsilon^{-1}$. This is the same advanced preconditioner used by MPB
and is the more complex of the two.

For TM, the structure is simpler. Because $A_{\mathrm{TM}}$ is already
diagonal in the plane-wave basis, its inverse is trivially available
and serves as an effective preconditioner. The result is that TM
consistently achieves lower condition numbers and converges in fewer
iterations than TE, a systematic advantage rooted in operator
structure.

Tuning the preconditioner accounted for roughly a quarter of the total
development time. The MPB developers describe the same challenge:
finding a good preconditioner ``is something of a \emph{black art} with
lots of trial and error''\footnote{From the MPB developer documentation,
\href{https://mpb.readthedocs.io/en/latest/Developer_Information/\#the-mathematics-of-mpb}{Developer
Information: The Mathematics of MPB}.}~\cite{mpb}. Because Blaze2D's
internal operator layout differs from MPB's, the preconditioner could
not be ported directly; it had to be re-derived and re-tuned against
condition-number measurements at each stage.

\section{The LOBPCG Solver}\label{sec:lobpcg}

At the heart of Blaze2D lies the LOBPCG algorithm~\cite{knyazev2001}.
At each step, the solver maintains three subspaces: the current
eigenvector candidates $X$ ($N_{\mathrm{band}}$ vectors), the search
directions $P$ from the previous step, and the preconditioned
residuals $W$. These are assembled into a trial basis $[X, P, W]$, and
a small dense eigenproblem (the Rayleigh--Ritz projection) of dimension
$\dim(X) + \dim(P) + \dim(W)$ is solved exactly. The new $X$ is read
off from the lowest eigenpairs of this projected problem. The structure
of this search subspace is sketched in \cref{fig:lobpcg_illustration}.

\begin{figure}[t]
  \centering
  \includegraphics[width=\linewidth]{lobpcg_subspace}
  \caption{Structure of the LOBPCG search subspace. The trial basis
    $[X, P, W]$ defines a small, dense projected eigenproblem whose
    exact solution yields the updated eigenvector approximations.
    Rather than operating in the full $N_{\mathrm{pw}}$-dimensional
    Hilbert space, each iteration works in a subspace of dimension
    $\dim(X) + \dim(P) + \dim(W)$.}
  \label{fig:lobpcg_illustration}
\end{figure}

Several implementation details proved decisive for performance and
stability:

\paragraph{Subspace dimensions.}
Expanding the candidate space beyond the number of requested bands
accelerates convergence significantly: adding one or two extra trial
vectors to $X$ provides the optimizer with additional variational
freedom at modest cost. Conversely, shrinking or omitting either $P$
or $W$ degrades convergence severely. The best results were achieved
when all three subspaces had roughly comparable dimensions.

\paragraph{Orthogonalization.}
Near-degenerate eigenvectors cause classical Gram--Schmidt
orthogonalization to lose numerical rank, producing a singular
Rayleigh--Ritz problem. Blaze2D uses the SVQB procedure~\cite{stathopoulos2002},
which orthogonalizes via a singular-value decomposition of the Gram
matrix and is numerically stable even for tightly clustered
eigenvalues.

\paragraph{Convergence criterion.}
MPB monitors the residual norm $\|A x - \lambda B x\|$ and declares
convergence when it falls below a threshold (default $10^{-7}$).
Blaze2D instead monitors the relative change in eigenvalues between
iterations. In practice, eigenvalues stabilize well before the residual
reaches a tight tolerance, so eigenvalue-based convergence requires
substantially fewer iterations per solve. While residual convergence
carries a proven convergence-rate guarantee~\cite{knyazev2001},
eigenvalue convergence proved completely stable across all tested
configurations, including difficult cases with accidental degeneracies
at high-symmetry $\bm{k}$-points.

\paragraph{Warm-starting.}
The single largest lever for performance is also the simplest:
initializing each $\bm{k}$-point solve with the converged eigenvectors
of the previous $\bm{k}$-point. This adds no algorithmic complexity,
yet it typically reduces the iteration count by an order of magnitude
compared to a random initialization. MPB uses the same strategy.

\paragraph{Deflation strategy.}
Once a band converges, the solver must avoid rediscovering it. Because
distinct Bloch modes are mathematically orthogonal, we can explicitly
project all new search directions away from the already-converged
states, a technique known as deflation. While solvers can
\emph{hard-lock} these converged bands by removing them from the active
problem entirely, Blaze2D relies on \emph{soft-locking}. Converged
bands remain part of the small Rayleigh--Ritz projection, but the
solver stops calculating new, computationally expensive search
directions for them. This halts their active optimization, yet allows
the converged states to mix with the active subspace to maintain
strict orthogonality as the remaining bands settle. In practice,
soft-locking proved far more numerically stable than hard-locking,
especially near accidental degeneracies where bands must be handled
carefully as a collective cluster rather than frozen one by one.

\section{Mixed Precision}\label{sec:mixed_precision}

Modern eigensolvers are not compute-bound but memory-bandwidth-bound.
Profiling both MPB and Blaze2D shows that actual floating-point
arithmetic accounts for as little as 10\% of CPU time; the remainder
is spent moving data between memory and cache. This observation
motivates a mixed-precision strategy: store all large arrays
(eigenvectors, operator matrices, intermediate products) in single
precision (\texttt{f32}) to halve the memory traffic, but promote to
double precision (\texttt{f64}) on the fly for numerically sensitive
operations (inner products, Rayleigh--Ritz projections, and
orthogonalization steps) before rounding the result back to
\texttt{f32} for storage.

The natural objection is: if the result is rounded back to \texttt{f32},
what is gained by computing in \texttt{f64}? The answer is that certain
intermediate accumulations require \texttt{f64} for numerical stability
(a dot product of $10^{4}$ terms, for instance, can lose several
significant digits in \texttt{f32}), while the \emph{result} of that
dot product (a single scalar or a small matrix) is representable in
\texttt{f32} without meaningful loss. The precision is needed
transiently, not persistently.

This mixed-precision approach yields a nearly clean $2\times$ speedup
with no significant loss in eigenvalue accuracy. The same conclusion
was reached independently in a 2023 study of mixed-precision LOBPCG for
electronic-structure calculations~\cite{woo2023}, a paper discovered
only after the approach had already been implemented in Blaze2D. The
strategy itself is standard in modern data science and
high-performance computing.

\section{Validation}\label{sec:validation}

Blaze2D was validated against MPB across a representative set of
two-dimensional photonic crystals by comparing full band diagrams,
$|\bm{u}|^{2}$ mode profiles, and convergence behavior. The validation
process was the most time-consuming part of the project. MPB uses
several nonstandard coordinate conventions, most notably a
half-pixel-shifted grid center, which produced false discrepancies
that had to be identified and accounted for one by one.

The validation narrative follows the same logic as the solver itself.
First, the global spectral output must agree. Second, the eigenvectors
must represent the same physical modes. Third, both solvers must
converge in the same way as the discretization is refined. The
following subsections address these three checks in turn.

\subsection{Band-Diagram Agreement}\label{sec:band_agreement}

The primary metric is the relative eigenvalue error across all
$\bm{k}$-points and bands. Bloch-function phases were not compared
directly, because MPB's phase conventions differ from Blaze2D's in
ways that scramble the complex phases without affecting physically
observable quantities. However, the squared modulus
$|u_{n\bm{k}}(\bm{r})|^{2}$ was compared and found to agree, and
improvements in mode-profile agreement tracked improvements in
eigenvalue accuracy, confirming that the band diagram is a sufficient
validation metric.

One systematic difference remains: when tracking bands across avoided
crossings, MPB tends to follow a given band adiabatically to higher
eigenvalues, while Blaze2D reports the true $N$ lowest eigenvalues at
each $\bm{k}$-point and does not track crossings. For isolated bands
away from topmost-band crossings, both solvers agree to high
precision. This behavior is visible in \cref{fig:band_diagram_agreement}.

\begin{figure}[t]
  \centering
  \includegraphics[width=\linewidth]{band_diagrams}
  \caption{Band-diagram comparison for a square lattice
    ($\epsilon = 8.9$ rods, $r = 0.2a$), showing the 10 lowest bands.
    \textit{Left}: TM polarization; \textit{right}: TE polarization.
    Both solvers show near-perfect agreement, with small deviations
    visible only in the topmost TE band due to differing band-tracking
    behavior across avoided crossings.}
  \label{fig:band_diagram_agreement}
\end{figure}

To quantify the agreement beyond the visual overlap of
\cref{fig:band_diagram_agreement}, \cref{fig:eigenvalue_comparison}
shows the RMS relative eigenvalue difference between the two solvers
for the lowest eight bands, evaluated over 61 $\bm{k}$-points along
$\Gamma \to X \to M \to \Gamma$, as a function of resolution. At
$N = 128$, relative differences reach $\approx 2 \times 10^{-5}$ on
average, with per-band maxima below $10^{-4}$. This is consistent with
Blaze2D's internal eigenvalue convergence tolerance of $10^{-5}$ in
mixed-precision mode; the residual plateau beyond $N \approx 48$
reflects that tolerance floor rather than a systematic discrepancy
between the two solvers.

\begin{figure}[t]
  \centering
  \includegraphics[width=\linewidth]{eigenvalue_comparison}
  \caption{RMS relative eigenvalue difference between Blaze2D and MPB
    for the eight lowest bands of a square lattice ($\epsilon = 8.9$
    rods, $r = 0.2a$). TM (left) and TE (right) polarizations.}
  \label{fig:eigenvalue_comparison}
\end{figure}

Taken together, \cref{fig:band_diagram_agreement} and
\cref{fig:eigenvalue_comparison} establish a precise agreement: at the
resolutions relevant for the registry sweep, Blaze2D reproduces MPB's
eigenvalues to within the mixed-precision solver tolerance.

\subsection{Mode Profiles}\label{sec:mode_profiles}

Eigenvalue agreement alone is not sufficient: the envelope pipeline
ultimately processes eigenvectors and matrix elements built from them.
For that reason, representative real-space mode profiles were compared
as well. \Cref{fig:u2_comparison} shows the squared Bloch amplitude
$|\bm{u}|^{2}$ for several bands and $\bm{k}$-points.

\begin{figure}[t]
  \centering
  \includegraphics[width=\linewidth]{u2_mode_comparison}
  \caption{Comparison of $|\bm{u}|^{2}$ mode profiles from Blaze2D and
    MPB for a representative set of bands, $\bm{k}$-points and a
    resolution of 64\,px/$a$. The spatial intensity distributions
    agree well for non-degenerate frequencies; at degenerate
    frequencies, the modes may be superposed or swapped, leading to
    expected visual discrepancies.}
  \label{fig:u2_comparison}
\end{figure}

The agreement is strong wherever the eigenvalue is non-degenerate:
nodal structure, localization, and symmetry match visually across the
unit cell. The only clearly visible mismatches occur at or near
degeneracies, where any orthonormal basis of the degenerate subspace
is equally valid. In those cases MPB and Blaze2D may return rotated or
swapped representatives of the same physical subspace, which changes
the plotted pattern without indicating a physical disagreement. This
is precisely why the comparison is carried out on $|\bm{u}|^{2}$
rather than on the complex Bloch phase itself.

Despite the strong agreement, small yet visual mismatches for bands 2
and 3 can be seen in \cref{fig:u2_comparison}. The spatial modes do
not appear perfectly symmetric; these are deviations from MPB that
have not been investigated or quantified here. Because the envelope
pipeline is built from Bloch-function data, these deviations should be
kept in mind when judging how reliable downstream envelope results
are.

\subsection{Convergence with Resolution}\label{sec:convergence}

Agreement at a single resolution could still be accidental. A stronger
test is whether both solvers approach the continuum limit at the same
rate as the plane-wave cutoff is increased. That comparison is shown
in \cref{fig:convergence_study}.

\begin{figure}[t]
  \centering
  \includegraphics[width=\linewidth]{convergence_study}
  \caption{Eigenvalue convergence for a square lattice
    ($\epsilon = 8.9$ rods, $r = 0.2a$). Solid lines: RMS deviation
    $\Delta\omega_{\mathrm{RMS}}$ of the four lowest bands along
    $\Gamma \to X \to M \to \Gamma$; dashed lines: worst-case
    deviation $\max|\Delta\omega|$. All deviations are measured
    against a higher resolution $N = 192$ reference. Both solvers
    converge at comparable rates.}
  \label{fig:convergence_study}
\end{figure}

The two curves track each other closely in both polarizations, for
both RMS and worst-case errors. This matters because it rules out a
hidden cancellation in the direct MPB comparison: Blaze2D does not
merely match one reference calculation, but exhibits the same
resolution dependence as an established solver. For the purposes of
the envelope pipeline, this is the relevant notion of correctness,
since all extracted matrix elements are ultimately built on these
converged Bloch states.

\section{Benchmarks}\label{sec:benchmarks}

Accuracy alone would not justify a new solver. The point of Blaze2D is
that it must also execute the dense registry sweep motivating it in
reasonable time. The runtime comparison is shown in
\cref{fig:blaze2d_benchmarks}, while the iteration counts,
per-iteration costs, and memory footprint are broken down in
\cref{fig:blaze2d_iterations_memory}.

\begin{figure}[t]
  \centering
  \includegraphics[width=\linewidth]{benchmarks}
  \caption{Single-core and multi-core wall-clock comparison of MPB and
    Blaze2D in mixed and full precision.}
  \label{fig:blaze2d_benchmarks}
\end{figure}

\begin{figure}[t]
  \centering
  \includegraphics[width=\linewidth]{iterations_memory}
  \caption{Iteration count per $\bm{k}$-point, time per iteration,
    and peak memory footprint for MPB vs.\ Blaze2D.}
  \label{fig:blaze2d_iterations_memory}
\end{figure}

A few observations from these benchmarks deserve comment. The TM
polarization converges in fewer iterations than TE for both solvers,
consistent with the simpler preconditioner structure discussed above.
Mixed precision does not increase the iteration count; the eigenvalue
trajectory is indistinguishable from the full-precision run. The
speedup seen in \cref{fig:blaze2d_benchmarks} therefore comes almost
entirely from reduced memory traffic rather than from altered
convergence behavior.

The decomposition in \cref{fig:blaze2d_iterations_memory} makes this
interpretation more concrete. Time per iteration falls under mixed
precision, while memory consumption scales predictably with
$N_{\mathrm{pw}}^{2}$ for the stored matrices and
$N_{\mathrm{pw}} \times N_{\mathrm{band}}$ for the eigenvector block,
with \texttt{f32} storage halving both contributions relative to
\texttt{f64}. For the registry sweep, this reduction matters twice: it
shortens each individual solve and increases the number of
configurations that can be processed concurrently.

\section{From Solver to Pipeline}\label{sec:pipeline}

With the eigensolver validated, the remaining task is to extract the
matrix elements required by the envelope effective Hamiltonian and
export them to the Python layer where the moir\'e-scale problem is
assembled and solved.

The extraction module computes, for each crystal configuration
$\bm{\delta}$ on the moir\'e grid, the local eigenvalues
$\lambda_{n}(\bm{\delta}, \bm{k}_{0})$, the Bloch functions
$u_{n\bm{k}_{0}}(\bm{r}; \bm{\delta})$, and the derived operator
matrix elements (Berry connections, effective-mass tensors, and the
coupling terms entering the envelope effective Hamiltonian). These are
written as structured arrays and exposed through Python bindings
generated via Rust's foreign-function interface.

Basic sanity checks on the extracted data (Hermiticity of the
effective Hamiltonian blocks, expected magnitudes of the
Berry-connection components, correct symmetry under lattice
operations) confirm that the plumbing between solver and pipeline
introduces no artifacts. At this point the narrative closes back onto
the envelope construction: the effective theory has been reduced to
computable local data, and Blaze2D is the tool that provides it.
